% Part: lambda-calculus
% Chapter: syntax
% Section: unique-readability

\documentclass[../../../include/open-logic-section]{subfiles}

\begin{document}

\olfileid{lam}{syn}{unq}
\olsection{وحدانية القراءة}

هل يمكن أن تتعدد طرق تكوين الحد الواحد؟ سنثبت أنها لا تتعدد:
فلكل حد لامبدا وجه واحد في البناء والتأويل. ومرادنا بـ\emph{التكوين}
فيما يلي بناء الحد بتطبيق قواعد \olref[trm]{defn:term} مرة أو مرات.

\begin{lem}\ollabel{lem:term-start}
لا يخلو أول الحد من أن يكون متغيرًا أو قوسًا.
\end{lem}

\begin{proof}
إنما يكون الشيء حدًا إذا بني بقواعد~\olref[trm]{defn:term}.
فما تبنيه \olref[trm]{defn:term-var} متغير، وما تبنيه
\olref[trm]{defn:term-abs} أو~\olref[trm]{defn:term-app} أوله قوس.
\end{proof}

\begin{lem}\ollabel{lem:app-start}
أول ناتج التطبيق قوسان، أو قوس يتلوه متغير.
\end{lem}

\begin{proof}
ليكن~$M$ حاصل تطبيق؛ فصورته~$(PQ)$، وأوله إذن قوس.
ثم إن~$P$ حد، فأوله، بموجب~\olref{lem:term-start}، قوس أو متغير.
\end{proof}

\begin{lem}\ollabel{lem:initial}
كل جزء ابتدائي حقيقي من حد يمتنع أن يكون حدًا.
\end{lem}

\begin{prob}
أقم برهان~\olref[lam][syn][unq]{lem:initial} مستعملًا الاستقراء على طول الحد.
\end{prob}

\begin{prop}[وحدانية القراءة] \ollabel{prop:unq}
تكوين كل حد وحيد؛ فإن انتهى تكوين إلى حد~$M$، امتنع أن ينتهي إليه تكوين آخر.
\end{prop}

\begin{proof}
  نستقرئ تكوين الحدود، ونميز الحالات الثلاث:
  \begin{enumerate}
    \item إذا كان~$M$ على صورة~$x$، حيث~$x$ متغير،
      لم يدخل التجريد ولا التطبيق في بنائه؛ إذ يبدأ ناتج كل منهما بقوس.
      فلا سبيل إلى بناء~$M$ إلا أن يكون تكوين~$M$ خطوة واحدة من
      \olref[trm]{defn:term}\olref[trm]{defn:term-var}.
    \item إذا كان~$M$ على صورة~$(\lambd[x][N])$، حيث~$x$ متغير و$N$ حد،
      امتنع أن تبنيه القاعدة
      \olref[trm]{defn:term}\olref[trm]{defn:term-var}؛ لأنه ليس متغيرًا مفردًا.
      ويمتنع أيضًا أن يكون حاصل تطبيق، بموجب~\olref{lem:app-start}.
      فليس~$M$ إلا حاصل تجريد~$N$، وتكوين~$N$ وحيد بفرض الاستقراء.
    \item إذا كان~$M$ على صورة~$(PQ)$، حيث~$P$ و$Q$ حدان،
      امتنع بناؤه بالقاعدة
      \olref[trm]{defn:term}\olref[trm]{defn:term-var}، لأن أوله قوس.
      وتقتضي~\olref{lem:term-start} أن~$P$ لا يبدأ بـ$\lambd$؛
      فلا يكون~$(PQ)$ حاصل تجريد. فلنفرض أن~$M$ بني بتطبيق آخر،
      فكانت صورته أيضًا~$(P'Q')$. إذا كان~$P \neq
      P'$، لزم أن يكون أحد~$P$ و$P'$
      جزءًا ابتدائيًا حقيقيًا من الآخر، وهو ممتنع بموجب~\olref{lem:initial}.
      إذن~$P=P'$، ومنه~$Q=Q'$. وبذلك يتعين~$P$ و$Q$،
      ثم ينفرد كل من~$P$ و$Q$ بتكوينه، بمقتضى فرض الاستقراء.
  \end{enumerate}
\end{proof}

وللقضية نفسها عبارة أيسر في القراءة:
\begin{prop}
  ينحصر الحد~$M$ في إحدى الصور التالية:
  \begin{enumerate}
    \item $x$، وفيها~$x$ متغير يتعين من~$M$ على وجه واحد.
    \item $(\lambd[x][N])$، وفيها~$x$ متغير و$N$ حد آخر، يتعين كلاهما من
      $M$ على وجه واحد.
    \item $(PQ)$، وفيها~$P$ و$Q$ حدان يتعينان من~$M$ على وجه واحد.
  \end{enumerate}
\end{prop}

\end{document}

